\section{Síntesis y ampliación}
\begin{table}[H]
\centering
\begin{tabular}{p{4cm}p{6cm}p{5cm}}
\toprule
Módulo & Aprendizaje clave & Siguiente paso \\
\midrule
Formalización matemática & Lean/Mathlib ofrece un deterministic environment y una gobernanza colaborativa & Incorporar el PR/gen replay log de Mathlib al experiment pipeline \\
Aprendizaje por refuerzo & El volante search-feedback-learning de AlphaZero es la base de diseño de AlphaProof & Continuar reforzando la search-policy distillation y el entropy gate \\
Arquitectura de AlphaProof & Formalizer + search + verifier + runner forman un pipeline controlable & Reforzar el mecanismo de activación de chunk gating y enriquecer el trace artifact \\
Evaluación y fallos & Las métricas multinivel y el failure curriculum convierten el failure en datos de entrenamiento & Ampliar el conjunto de métricas del dashboard y activar automáticamente el failure curriculum \\
Práctica de ingeniería & Proof search scheduling + trace logging + governance closing loop & Establecer el drift threshold y tratar cada alert como un retraining seed \\
Compartición de conocimiento & El log, el newsletter y el transcript forman un network de reproducibilidad & Hacer que el newsletter sea el archival record de los highlights de QA y de governance \\
Enseñanza del sistema & La arquitectura de tres capas de Pipeline/observability/governance constituye la base de la confianza & Incorporar la differential governance policy al newsletter \& la automation \\
\bottomrule
\end{tabular}
\caption{Módulos técnicos y plan de acción de la Lecture 05}
\end{table}

\subsection{Lecturas adicionales}
\begin{itemize}
    \item Hubert et al., ``AlphaProof: RL Meets Formal Mathematics,'' 2024.
    \item Documentación oficial de Lean y el README de Mathlib (GitHub).
    \item Artículos relacionados con AlphaZero / AlphaTensor.
    \item Archive público de silver proofs del IMO 2024.
    \item Blog de Terence Tao sobre la colaboración en formalización.
    \item Workshop notes de «Governance for autonomous agents».
\end{itemize}

\subsection{Resumen rápido de términos clave}
\begin{itemize}
    \item Chunk gating: selector de estrategia de token/overlap/latency dentro del proof search.
    \item Drift gate: la caída de entropy activa el fallback + review.
    \item Trace log: la tactic sequence + human comment de cada failed proof.
    \item Knowledge pipeline: el triángulo de log + newsletter + transcript.
\end{itemize}
\begin{knowledgebox}{El papel del resumen rápido de términos clave}
Usa una terminología unificada para ayudar a los nuevos integrantes a familiarizarse rápidamente con el operational playbook de AlphaProof, y se combina con el governance board para formar un ciclo cerrado de retroalimentación.
\end{knowledgebox}

\subsection{Resumen de la sección}
AlphaProof es un sistema que entrelaza la formalización matemática, el aprendizaje por refuerzo y una práctica de ingeniería gobernable: Lean ofrece el sandbox, RL ofrece el volante de search-policy, la gobernanza ofrece trace/alert/knowledge, y el newsletter mantiene sincronizado a todo el equipo.

\end{document}
